\section{連続性}
本節は位相空間論パートの連続性の稿も参照しておくとよい.

\subsection{連続性・一様連続性}



\begin{prop} (連続関数まとめ)
\ben
\item (\tf{中間値の定理})$I$を$\mb{R}$の区間として, 連続関数$f(x):I\to \mb{R}$, $a,b\in I$に対し, $f(a)$と$f(b)$の間の任意の実数値を$f(x)$が取る.
\item (\tf{最大値の原理}) $I$を有界閉区間とするとき, $f(x)$には最大値と最小値が存在する.
\item (\tf{平均値の定理})$f(x)$が閉区間$[a,b]$で連続, 開区間$(a,b)$で微分可能のとき, $f(b)-f(a) = f'(c)(b-a)$となる$c\in [a,b]$が存在する.
\item (\tf{超重要})\tf{コンパクト距離空間上の連続関数は一様連続である.}
\item (\tf{Borzano-Weierstrass})
閉区間$[0,1]\subset \mb{R}$は(点列)コンパクト集合である.
\een
\end{prop}
\prf
(4): $X$をコンパクト距離空間とする. 距離空間ではコンパクト性と点列コンパクト性が同値である. {\large \tf{主張を背理法で示す.}}\footnote{背理法以外でやると泥沼になる(経験談).} $f:X\to \mb{R}$を連続関数とせよ. もし一様連続でないなら, ある$\epsilon>0$について次が正しくなる.
\[\forall n\in \mb{N}, \exists x_n,y_n \in X, \mr{s.t.} |x_n - y_n| < \dfrac{1}{n} かつ |f(x_n) - f(y_n)| \geq \epsilon \]
(点列)コンパクト性より, $\set{x_n}$から収束部分列$\set{x_{v(n)}}$が取れる. この収束先を$\alpha$とすると, $y_{v(n)}$も$\alpha$に収束する. そして, $f$の連続性により
\[f(x_{v(n)}) - f(y_{v(n)}) \to f(\alpha) - f(\alpha) = 0\]
だが, これは誤差$\epsilon$以上という条件に反する.
\qed

\begin{egprob}
\ben
\item 連続関数$f:S^1 \to \mb{R}$は全射にも単射にもなり得ないことを示せ.
\item 周期関数には最大値と最小値が存在することを示せ.
\item $f(x)$が$a$を含む区間で$C^n$級であるとき, $R_n = f^{(n)}(a)(x-a)^n/n! + \epsilon(x) (x-a)^n$とおく. ただし$\epsilon(x)$は
\[f(x) = f(a) + f'(a)(x-a) + \dots + \dfrac{f^{(n-1)}(a)}{(n-1)!}(x-a)^{n-1} + R_n(x)\]
を満たすとする. このとき$x\to a $で$\epsilon(x) \to 0$ であることを示せ.
\item 連続関数$f:\mb{R}\to \mb{R}$で一様連続でないものは? $f$が有界の場合は?
\item $\mb{R}$のCauchy列は収束列であることを示せ.
\een
\end{egprob}
\begin{egans}
(1): 全射にならないのは最大値と最小値があるから. 単射にならないのは, 最小点と最大点を結ぶ経路が二つ取れて, その経路の間の値を取るので, 少なくとも2回$f$の取る値が存在するため. \\
(2): 周期を$T$とすると, $[0,T]$上で値が決まるから最大値原理を使えばよい.  \\
(3): $R^{(k)}(a) = 0$ ($k=0,1,\dots, n-1$)に注意.
$R_n$は$C^n$級なので, $R_n(x) - R_n(a) = (x-a)R_n'(c_x)$を満たす$c_{x,1}$がある.
$R_n'(c_{x,1})=R_n'(c_{x,1}) - R_n'(a) = (x-a)R_n''(c_{x,2})$と書ける. これを続け, $R_n(x) = (x-a)^n R_n^{(n)}(c_{x,n})$となる$a<c_{x,n}<x$がある.
よって
\[R^{(n)}(c_{x,n}) = f^{(n)}(a) + \epsilon(x) \quad (x\neq a)\]
であるが, $R_n$を$f(x) - f(a) - \dots - \frac{f^{(n-1)}(a)}{(n-1)!}(x-a)^{n-1}$とみると $R^{(n)}(c_{x,n}) = f^{(n)}(c_{x,n})$である. $f^{(n)}$は連続だから, $\epsilon(x) = R^{(n)}(c_{x,n}) - f^{(n)}(a)\to 0$である. \qed
\\
(4): $f(x) = \sin{x^2}$とせよ. \\
(5): Cauchy列$\set{a_n}$を取る. $a_n$は有界列であるから有界閉区間上の点列だと思える. \tf{このとき収束部分列が取れる}が, Cauchy列が収束部分列を持てば元々収束列にもなる. 実際, $\epsilon > 0$を任意に取り, $a_{v(n)}\to \alpha$なる収束部分列があるとき, 十分大きい$N$について, $n\in \mb{N}$が$n,v(n) > N$を満たすときに$|a_{n} - a_{v(n)}| < \epsilon$である.
\[|a_N - \alpha| < |a_n - a_{v(n)}| + |a_{v(n)} - \alpha| < 2\epsilon \]
とできるから.
\end{egans}


次は自分で考えてみてほしい.
\begin{prac} (仮想面接問題)
$\mb{R}$上の可積分関数は有界ですか?
\end{prac}

この問題に関連して次のような問題がある.
\begin{egprob} (2021東北大共通, 改題)
$f:\mb{R}\to\mb{R}$は$\mb{R}$上微分可能であり, $f'$は$\mb{R}$で有界でかつ, $\int_{-\infty}^{\infty} |f(t)|\, dt$は収束するとする. このとき, $f$は$\mb{R}$で一様連続であることを示せ. また, $f(t)$は$t\to \infty$とするとき0に収束することを示せ.
\end{egprob}
\begin{egans}
\chatgpthint{$f'$の有界性から一様連続性を示し、反証法で積分可能性と矛盾させてください。}
\end{egans}



演習問題でときに見かける定理と,その応用例を記しておく.
\begin{theo} (\tf{Diniの定理}) \label{theo:Dini_theorem}
閉区間$I=[a,b]$上の連続な関数列$\{ f_n(x)\}$が$I$上単調に増加しながら$f(x)$に収束するとする.このとき, $f(x)$が$I$上連続であったなら, $f_n(x)$は$I$上一様に$f_n(x)$に収束する.単調減少の場合も同様である.
\end{theo}
\begin{egprob} (\cite{rikou}演習5-27, 改)
$f_n(x) = (1-x^2/n)^{n}$ ($0\leq x\leq \sqrt{n}$), $=0$ ($\sqrt{n}\leq x$)とする.$f_n(x)$は単調に増加しながら$e^{-x^2}$に収束することを確かめ, 次の極限を求めよ.
\[\disp\lim_{n\to \infty} \int_{0}^{\infty} f_n(x)\, dx\]
\end{egprob}
\begin{egans}
\chatgpthint{各点収束を確認し、優収束定理で極限関数を積分してください。}
\end{egans}


\subsection{中間値の定理・最大値の原理}
連続性に関する代表的な二つの定理です. 主張はよく知られているため省略します.
\begin{egprob}
連続写像$f:S^1\to \mb{R}$は全射にも単射にもならないことを示せ.
\end{egprob}
\begin{egans}
$S^1$はコンパクトなので最大値の原理より$f$は最大値$M$と最小値$m$を持つ. よって全射になり得ない. $f$が単射だとせよ. $f(x) = M$, $f(y) = m$だとする. $M=m$なら定数なので単射でないから$M>m$とする. $x,y$は異なる点であるが, $x$と$y$を結ぶ円弧が二つある. 中間値の定理より, それぞれの円弧上で, $\frac{m+M}{2}$という値を取る点があるので, 単射になり得ない.
\end{egans}







\subsection{演習問題}

\begin{prob} (H31基礎科目6) \label{prob:H31kiso6}
$\mb{R}^2$上の実数値連続関数$f$についての次の条件を考える,
\begin{center}
    ($\ast$) 任意の$R>0$に対して次の集合は有界である.
    \[\set{(x,y)\in \mb{R}^2 \mid |f(x,y)| \leq R}\]
\end{center}
\ben
\item 条件$(\ast)$を満たす連続関数$f$の例を与え, それが$(\ast)$をみたすことを示せ.
\item 連続関数$f$が条件$(\ast)$をみたすとき, 次のいずれかが成り立つことを示せ.
\ben
\item $f$は最大値を持つが, 最小値は持たない.
\item $f$は最小値を持つが, 最大値は持たない.

\een
\een
\end{prob}

\subsection{Hint・Answer}
{\small
\begin{probans}
最大値と最小値を両方持つことはない. 実際,そのとき $|f|$が有界であり$\set{|f| < R} = \mb{R}^2$となる$R$が取れるため. 最大値と最小値を両方持たないとする(\tf{注意}: これはいくらでも大きい・小さい値を取るという条件にはならない！). \par
$R>0$に対し$K_R = \set{\sqrt{x^2+y^2}\leq R}$とおく. 条件より$f$は定数関数ではありえないので, $f$はある2点で異なる値を取る. その2点を含む$K_r$が取れ, $f|_{K_r}$は最大値$M_r$と最小値$N_r$を持ち, 取り方から$M_r>N_r$である. $N_r< c< M_r$なる実数$c$を任意に一つ取って固定する. $\set{f=c}$は有界集合なので, 十分大きい$R$では$\set{f=c} \subset K_R$となる. $R>r$としてもよく, $f|_{K_R}$の最大値$M_R$, 最小値$N_R$は$M_R \geq M_r > c > N_r \geq N_{R}$を満たす. ところが$f$は最大値と最小値を持たないので, $M_R,N_R$は$f$の最大値,
最小値にはならない. よって, ある$P_{R},Q_{R}\in \mb{R}^2$が存在して$f(P_{R}) > M_R$, $f(Q_{R}) < N_R$となる. $M_R,N_R$は$K_R$上の最大値・最小値なので$P_R,Q_R\notin K_R$である. $K_R$は円盤なので, $P_R,Q_R$を結ぶ曲線$\gamma_{R}:[0,1]\to \mb{R}^2$で$K_R$と交わらないものが存在する. $\gamma_{R}(0) = P_R$, $\gamma_{R}(1) = Q_R$とすると $f\circ \gamma_{R}:[0,1]\to \mb{R}$が$0\mapsto M_R, 1\mapsto N_R$なる連続関数で, 中間値の定理よりある$t_R\in [0,1]$に対して$f\circ \gamma_R(t_R) = c$を満たさなければならない. $\set{f=c} \subset K_R$より$\gamma_R(t_R)\in K_R$であるが, これは$\gamma_R$が$K_R$と交わらないという仮定に反する. \qed
\end{probans}
}








\clearpage
%積分パート
%積分の理論, 重積分計算問題, 広義積分の理論
